% --- Archon physics formalization source begin ---
% archon:physics
% archon:covers PhyXMiniProblems/problem_phyx_mini_0515.lean
% archon:source-report reports/phyx_mini/problem_phyx_mini_0515.source.json

\chapter{Physics problem phyx\_mini\_0515}
\label{ch:PhyXMiniProblems_problem_phyx_mini_0515}

\paragraph{Problem source.}
\#\# Physical scenario

The graph in figure shows the stopping potential as a function of the frequency of the incident light falling on a metal surface.

\#\# Question

Find the photoelectric work function for this metal.

\#\# Answer choices

- A: A: \textbackslash{}( 5.6 \textbackslash{}text\{ eV\} \textbackslash{})

- B: B: \textbackslash{}( 2.9 \textbackslash{}text\{ eV\} \textbackslash{})

- C: C: \textbackslash{}( 3.2 \textbackslash{}text\{ eV\} \textbackslash{})

- D: D: \textbackslash{}( 4.8 \textbackslash{}text\{ eV\} \textbackslash{})

\#\# Figure caption (auxiliary; use the image as primary evidence)

The image is a graph with the title "Stopping potential (V)." The graph is a plot with the x-axis labeled from 0 to 3.0 in increments of 0.50 and the y-axis labeled from 0.0 to 7.0 in increments of 1.0.

- **X-axis:** It likely represents an independent variable but is not labeled with a specific quantity or unit.
- **Y-axis:** Labeled as "Stopping potential (V)," suggesting it represents potential in volts.
- **Line:** A straight line runs through the graph, starting from a point near (0, 0) and extending upwards, indicating a positive linear relationship.
- **Grid:** A grid is present, aiding in reading values accurately.

There's no additional text or objects beyond the graph components.

\#\# Dataset metadata

Quantum Phenomena; Physical Model Grounding Reasoning, Multi-Formula Reasoning

\paragraph{Recorded answer/context.}
D: D: \textbackslash{}( 4.8 \textbackslash{}text\{ eV\} \textbackslash{})

\paragraph{Figure/image path.}
/root/proposal\_for\_physic/science-mango/phyx\_mini\_run/phyx\_data/test\_image/515.png

\paragraph{Formalization target.}
create a compiling Lean file with sorry bodies at `PhyXMiniProblems/problem\_phyx\_mini\_0515.lean`. The Lean declarations must preserve the physical quantities, dimensions or dimensional roles, figure labels, governing-law hypotheses, and final relation expressed by this problem.
Use Mathlib/Physlib names found through LeanExplore where available. If a physics API is missing, introduce faithful local abstractions rather than scalar placeholder aliases.

\begin{theorem}[Physics formalization target]
\label{thm:physics:phyx_mini_0515:target}
\lean{PhyXMiniProblems.ProblemPhyXMini0515.photoelectric_workFunction_matches_choiceD}
\uses{def:physics:phyx-mini-0515:phyxminiproblems-problemphyxmini0515-frequencyinhertz, def:physics:phyx-mini-0515:phyxminiproblems-problemphyxmini0515-actioninjouleseconds, def:physics:phyx-mini-0515:phyxminiproblems-problemphyxmini0515-energyinjoules, def:physics:phyx-mini-0515:phyxminiproblems-problemphyxmini0515-energyinelectronvolts, def:physics:phyx-mini-0515:phyxminiproblems-problemphyxmini0515-photoelectricsetup, def:physics:phyx-mini-0515:phyxminiproblems-problemphyxmini0515-matchesphotoelectricscenario, def:physics:phyx-mini-0515:phyxminiproblems-problemphyxmini0515-matchessuppliedstoppingpotentialgraph, def:physics:phyx-mini-0515:phyxminiproblems-problemphyxmini0515-usesstandardphotoelectricreferencedata, def:physics:phyx-mini-0515:phyxminiproblems-problemphyxmini0515-hasphysicalphotoelectricparameters, def:physics:phyx-mini-0515:phyxminiproblems-problemphyxmini0515-satisfieseinsteinphotoelectriclaw, def:physics:phyx-mini-0515:phyxminiproblems-problemphyxmini0515-displayedworkfunctioninelectronvolts, def:physics:phyx-mini-0515:phyxminiproblems-problemphyxmini0515-isclosestdisplayedworkfunction}
The assigned autoformalize agent should translate this physics problem into Lean declarations in the covered file, with theorem and lemma proof bodies written as `by sorry`.
\end{theorem}
\begin{proof}
This is an autoformalization task, not a proof task. Produce statements that can later be proved by the physics prover without weakening the physical model.
\end{proof}

% --- Archon declaration topology begin ---
\section{Declaration topology}
\begin{definition}[Frequency Quantity]
  \label{def:physics:phyx-mini-0515:phyxminiproblems-problemphyxmini0515-frequencyquantity}
  \lean{PhyXMiniProblems.ProblemPhyXMini0515.FrequencyQuantity}
  A physical cyclic frequency, with inverse-time dimension
\end{definition}
\begin{proof}
  This is the declared model definition of Frequency Quantity.
\end{proof}
\begin{definition}[Charge Magnitude Quantity]
  \label{def:physics:phyx-mini-0515:phyxminiproblems-problemphyxmini0515-chargemagnitudequantity}
  \lean{PhyXMiniProblems.ProblemPhyXMini0515.ChargeMagnitudeQuantity}
  A physical electric-charge magnitude
\end{definition}
\begin{proof}
  This is the declared model definition of Charge Magnitude Quantity.
\end{proof}
\begin{definition}[potential Difference Dimension]
  \label{def:physics:phyx-mini-0515:phyxminiproblems-problemphyxmini0515-potentialdifferencedimension}
  \lean{PhyXMiniProblems.ProblemPhyXMini0515.potentialDifferenceDimension}
  The dimension M L² T⁻² C⁻¹ of an electric potential difference
\end{definition}
\begin{proof}
  This is the declared model definition of potential Difference Dimension.
\end{proof}
\begin{definition}[Potential Difference Quantity]
  \label{def:physics:phyx-mini-0515:phyxminiproblems-problemphyxmini0515-potentialdifferencequantity}
  \lean{PhyXMiniProblems.ProblemPhyXMini0515.PotentialDifferenceQuantity}
  \uses{def:physics:phyx-mini-0515:phyxminiproblems-problemphyxmini0515-potentialdifferencedimension}
  A physical stopping-potential difference
\end{definition}
\begin{proof}
  This is the declared model definition of Potential Difference Quantity.
\end{proof}
\begin{definition}[action Dimension]
  \label{def:physics:phyx-mini-0515:phyxminiproblems-problemphyxmini0515-actiondimension}
  \lean{PhyXMiniProblems.ProblemPhyXMini0515.actionDimension}
  The dimension M L² T⁻¹ of action
\end{definition}
\begin{proof}
  This is the declared model definition of action Dimension.
\end{proof}
\begin{definition}[Action Quantity]
  \label{def:physics:phyx-mini-0515:phyxminiproblems-problemphyxmini0515-actionquantity}
  \lean{PhyXMiniProblems.ProblemPhyXMini0515.ActionQuantity}
  \uses{def:physics:phyx-mini-0515:phyxminiproblems-problemphyxmini0515-actiondimension}
  A physical action, used here for the ordinary Planck constant h
\end{definition}
\begin{proof}
  This is the declared model definition of Action Quantity.
\end{proof}
\begin{definition}[frequency In Hertz]
  \label{def:physics:phyx-mini-0515:phyxminiproblems-problemphyxmini0515-frequencyinhertz}
  \lean{PhyXMiniProblems.ProblemPhyXMini0515.frequencyInHertz}
  \uses{def:physics:phyx-mini-0515:phyxminiproblems-problemphyxmini0515-frequencyquantity}
  SI readout of a frequency in hertz
\end{definition}
\begin{proof}
  This is the declared model definition of frequency In Hertz.
\end{proof}
\begin{definition}[charge In Coulombs]
  \label{def:physics:phyx-mini-0515:phyxminiproblems-problemphyxmini0515-chargeincoulombs}
  \lean{PhyXMiniProblems.ProblemPhyXMini0515.chargeInCoulombs}
  \uses{def:physics:phyx-mini-0515:phyxminiproblems-problemphyxmini0515-chargemagnitudequantity}
  SI readout of a charge magnitude in coulombs
\end{definition}
\begin{proof}
  This is the declared model definition of charge In Coulombs.
\end{proof}
\begin{definition}[potential Difference In Volts]
  \label{def:physics:phyx-mini-0515:phyxminiproblems-problemphyxmini0515-potentialdifferenceinvolts}
  \lean{PhyXMiniProblems.ProblemPhyXMini0515.potentialDifferenceInVolts}
  \uses{def:physics:phyx-mini-0515:phyxminiproblems-problemphyxmini0515-potentialdifferencequantity}
  SI readout of a potential difference in volts
\end{definition}
\begin{proof}
  This is the declared model definition of potential Difference In Volts.
\end{proof}
\begin{definition}[action In Joule Seconds]
  \label{def:physics:phyx-mini-0515:phyxminiproblems-problemphyxmini0515-actioninjouleseconds}
  \lean{PhyXMiniProblems.ProblemPhyXMini0515.actionInJouleSeconds}
  \uses{def:physics:phyx-mini-0515:phyxminiproblems-problemphyxmini0515-actionquantity}
  SI readout of an action in joule-seconds
\end{definition}
\begin{proof}
  This is the declared model definition of action In Joule Seconds.
\end{proof}
\begin{definition}[energy In Joules]
  \label{def:physics:phyx-mini-0515:phyxminiproblems-problemphyxmini0515-energyinjoules}
  \lean{PhyXMiniProblems.ProblemPhyXMini0515.energyInJoules}
  SI readout of a Physlib energy in joules
\end{definition}
\begin{proof}
  This is the declared model definition of energy In Joules.
\end{proof}
\begin{definition}[energy In Electron Volts]
  \label{def:physics:phyx-mini-0515:phyxminiproblems-problemphyxmini0515-energyinelectronvolts}
  \lean{PhyXMiniProblems.ProblemPhyXMini0515.energyInElectronVolts}
  \uses{def:physics:phyx-mini-0515:phyxminiproblems-problemphyxmini0515-energyinjoules}
  Read an energy in electron-volts. Physlib's DimEnergy.electronVolt is the dimensionful energy corresponding to one electron-volt, so this is a ratio of two energy readouts rather than a scalar alias for energy
\end{definition}
\begin{proof}
  This is the declared model definition of energy In Electron Volts.
\end{proof}
\begin{definition}[Plot Axis]
  \label{def:physics:phyx-mini-0515:phyxminiproblems-problemphyxmini0515-plotaxis}
  \lean{PhyXMiniProblems.ProblemPhyXMini0515.PlotAxis}
  The two axes shown in the supplied graph
\end{definition}
\begin{proof}
  This is the declared model definition of Plot Axis.
\end{proof}
\begin{definition}[Plot Axis Role]
  \label{def:physics:phyx-mini-0515:phyxminiproblems-problemphyxmini0515-plotaxisrole}
  \lean{PhyXMiniProblems.ProblemPhyXMini0515.PlotAxisRole}
  Physical meaning assigned to each graph axis
\end{definition}
\begin{proof}
  This is the declared model definition of Plot Axis Role.
\end{proof}
\begin{definition}[Stopping Potential Figure]
  \label{def:physics:phyx-mini-0515:phyxminiproblems-problemphyxmini0515-stoppingpotentialfigure}
  \lean{PhyXMiniProblems.ProblemPhyXMini0515.StoppingPotentialFigure}
  \uses{def:physics:phyx-mini-0515:phyxminiproblems-problemphyxmini0515-plotaxis, def:physics:phyx-mini-0515:phyxminiproblems-problemphyxmini0515-plotaxisrole}
  The directly visible features of the supplied graph. Coordinates are the dimensionless numbers printed on the axes. The horizontal physical scale is connected to hertz separately in MatchesSuppliedStoppingPotentialGraph
\end{definition}
\begin{proof}
  This is the declared model definition of Stopping Potential Figure.
\end{proof}
\begin{definition}[Photoelectric Setup]
  \label{def:physics:phyx-mini-0515:phyxminiproblems-problemphyxmini0515-photoelectricsetup}
  \lean{PhyXMiniProblems.ProblemPhyXMini0515.PhotoelectricSetup}
  \uses{def:physics:phyx-mini-0515:phyxminiproblems-problemphyxmini0515-frequencyquantity, def:physics:phyx-mini-0515:phyxminiproblems-problemphyxmini0515-chargemagnitudequantity, def:physics:phyx-mini-0515:phyxminiproblems-problemphyxmini0515-potentialdifferencequantity, def:physics:phyx-mini-0515:phyxminiproblems-problemphyxmini0515-actionquantity, def:physics:phyx-mini-0515:phyxminiproblems-problemphyxmini0515-stoppingpotentialfigure}
  The physical objects and unknown quantities in the photoelectric experiment. The work function is a field, not a definition in terms of an answer choice
\end{definition}
\begin{proof}
  This is the declared model definition of Photoelectric Setup.
\end{proof}
\begin{definition}[Matches Photoelectric Scenario]
  \label{def:physics:phyx-mini-0515:phyxminiproblems-problemphyxmini0515-matchesphotoelectricscenario}
  \lean{PhyXMiniProblems.ProblemPhyXMini0515.MatchesPhotoelectricScenario}
  \uses{def:physics:phyx-mini-0515:phyxminiproblems-problemphyxmini0515-photoelectricsetup}
  The qualitative physical scenario stated in the problem
\end{definition}
\begin{proof}
  This is the declared model definition of Matches Photoelectric Scenario.
\end{proof}
\begin{definition}[Matches Supplied Stopping Potential Graph]
  \label{def:physics:phyx-mini-0515:phyxminiproblems-problemphyxmini0515-matchessuppliedstoppingpotentialgraph}
  \lean{PhyXMiniProblems.ProblemPhyXMini0515.MatchesSuppliedStoppingPotentialGraph}
  \uses{def:physics:phyx-mini-0515:phyxminiproblems-problemphyxmini0515-frequencyinhertz, def:physics:phyx-mini-0515:phyxminiproblems-problemphyxmini0515-potentialdifferenceinvolts, def:physics:phyx-mini-0515:phyxminiproblems-problemphyxmini0515-photoelectricsetup}
  Axis labels, tick labels, qualitative trace shape, and the threshold intercept read from the primary image. The graph's prose identifies the unlabeled horizontal axis as frequency; its numerical coordinate is calibrated in units of 10\textasciicircum{}15 Hz. None of these fields states the work function or an answer choice
\end{definition}
\begin{proof}
  This is the declared model definition of Matches Supplied Stopping Potential Graph.
\end{proof}
\begin{definition}[Uses Standard Photoelectric Reference Data]
  \label{def:physics:phyx-mini-0515:phyxminiproblems-problemphyxmini0515-usesstandardphotoelectricreferencedata}
  \lean{PhyXMiniProblems.ProblemPhyXMini0515.UsesStandardPhotoelectricReferenceData}
  \uses{def:physics:phyx-mini-0515:phyxminiproblems-problemphyxmini0515-chargeincoulombs, def:physics:phyx-mini-0515:phyxminiproblems-problemphyxmini0515-actioninjouleseconds, def:physics:phyx-mini-0515:phyxminiproblems-problemphyxmini0515-photoelectricsetup}
  Standard SI calibration data. ordinaryPlanckAction uses the exact SI value of h, and the charge magnitude is the elementary charge. These reference values do not mention the unknown work function or the displayed answers
\end{definition}
\begin{proof}
  This is the declared model definition of Uses Standard Photoelectric Reference Data.
\end{proof}
\begin{definition}[Has Physical Photoelectric Parameters]
  \label{def:physics:phyx-mini-0515:phyxminiproblems-problemphyxmini0515-hasphysicalphotoelectricparameters}
  \lean{PhyXMiniProblems.ProblemPhyXMini0515.HasPhysicalPhotoelectricParameters}
  \uses{def:physics:phyx-mini-0515:phyxminiproblems-problemphyxmini0515-frequencyinhertz, def:physics:phyx-mini-0515:phyxminiproblems-problemphyxmini0515-chargeincoulombs, def:physics:phyx-mini-0515:phyxminiproblems-problemphyxmini0515-potentialdifferenceinvolts, def:physics:phyx-mini-0515:phyxminiproblems-problemphyxmini0515-actioninjouleseconds, def:physics:phyx-mini-0515:phyxminiproblems-problemphyxmini0515-energyinjoules, def:physics:phyx-mini-0515:phyxminiproblems-problemphyxmini0515-photoelectricsetup}
  Positivity conditions selecting a physically meaningful experiment
\end{definition}
\begin{proof}
  This is the declared model definition of Has Physical Photoelectric Parameters.
\end{proof}
\begin{definition}[Satisfies Einstein Photoelectric Law]
  \label{def:physics:phyx-mini-0515:phyxminiproblems-problemphyxmini0515-satisfieseinsteinphotoelectriclaw}
  \lean{PhyXMiniProblems.ProblemPhyXMini0515.SatisfiesEinsteinPhotoelectricLaw}
  \uses{def:physics:phyx-mini-0515:phyxminiproblems-problemphyxmini0515-frequencyinhertz, def:physics:phyx-mini-0515:phyxminiproblems-problemphyxmini0515-chargeincoulombs, def:physics:phyx-mini-0515:phyxminiproblems-problemphyxmini0515-potentialdifferenceinvolts, def:physics:phyx-mini-0515:phyxminiproblems-problemphyxmini0515-actioninjouleseconds, def:physics:phyx-mini-0515:phyxminiproblems-problemphyxmini0515-energyinjoules, def:physics:phyx-mini-0515:phyxminiproblems-problemphyxmini0515-photoelectricsetup}
  Einstein's photoelectric equation e V\_s = h f - φ for frequencies at or above threshold. It is a general governing relation between the stopping-potential trace and the independent work-function field; it does not supply the numerical answer requested by the problem
\end{definition}
\begin{proof}
  This is the declared model definition of Satisfies Einstein Photoelectric Law.
\end{proof}
\begin{definition}[Answer Choice]
  \label{def:physics:phyx-mini-0515:phyxminiproblems-problemphyxmini0515-answerchoice}
  \lean{PhyXMiniProblems.ProblemPhyXMini0515.AnswerChoice}
  Labels attached to the four displayed work-function choices
\end{definition}
\begin{proof}
  This is the declared model definition of Answer Choice.
\end{proof}
\begin{definition}[displayed Work Function In Electron Volts]
  \label{def:physics:phyx-mini-0515:phyxminiproblems-problemphyxmini0515-displayedworkfunctioninelectronvolts}
  \lean{PhyXMiniProblems.ProblemPhyXMini0515.displayedWorkFunctionInElectronVolts}
  \uses{def:physics:phyx-mini-0515:phyxminiproblems-problemphyxmini0515-answerchoice}
  The energy, in electron-volts, printed beside each answer label
\end{definition}
\begin{proof}
  This is the declared model definition of displayed Work Function In Electron Volts.
\end{proof}
\begin{definition}[Is Closest Displayed Work Function]
  \label{def:physics:phyx-mini-0515:phyxminiproblems-problemphyxmini0515-isclosestdisplayedworkfunction}
  \lean{PhyXMiniProblems.ProblemPhyXMini0515.IsClosestDisplayedWorkFunction}
  \uses{def:physics:phyx-mini-0515:phyxminiproblems-problemphyxmini0515-energyinelectronvolts, def:physics:phyx-mini-0515:phyxminiproblems-problemphyxmini0515-answerchoice, def:physics:phyx-mini-0515:phyxminiproblems-problemphyxmini0515-displayedworkfunctioninelectronvolts}
  A displayed choice is best when its printed electron-volt value is at least as close to the inferred work function as every other displayed value. This encodes the resolution of the graph-based multiple-choice question without assuming that the physical work function was defined by a choice
\end{definition}
\begin{proof}
  This is the declared model definition of Is Closest Displayed Work Function.
\end{proof}
\begin{lemma}[work Function eq planck Action mul threshold Frequency]
  \label{lem:physics:phyx-mini-0515:phyxminiproblems-problemphyxmini0515-workfunction-eq-planckaction-mul-thresholdfrequency}
  \lean{PhyXMiniProblems.ProblemPhyXMini0515.workFunction_eq_planckAction_mul_thresholdFrequency}
  \uses{def:physics:phyx-mini-0515:phyxminiproblems-problemphyxmini0515-frequencyinhertz, def:physics:phyx-mini-0515:phyxminiproblems-problemphyxmini0515-actioninjouleseconds, def:physics:phyx-mini-0515:phyxminiproblems-problemphyxmini0515-energyinjoules, def:physics:phyx-mini-0515:phyxminiproblems-problemphyxmini0515-photoelectricsetup, def:physics:phyx-mini-0515:phyxminiproblems-problemphyxmini0515-matchessuppliedstoppingpotentialgraph, def:physics:phyx-mini-0515:phyxminiproblems-problemphyxmini0515-satisfieseinsteinphotoelectriclaw}
  At the threshold, the zero stopping-potential readout and Einstein equation imply φ = h f₀. This is a derived physical relation, not a premise field
\end{lemma}
\begin{proof}
  The conclusion follows from the governing relations and figure readouts named in the statement of work Function eq planck Action mul threshold Frequency.
\end{proof}
% --- Archon declaration topology end ---
% --- Archon physics formalization source end ---
